The deep BSDE method~\cite{han2018solving} exploits the nonlinear Feynman-Kac
lemma: for a semilinear parabolic PDE, the solution $V(t,x_t)$ is the
$Y$-component of a backward stochastic differential equation (BSDE):
\begin{equation}
  dY_t = -f(X_t,Z_t)\,dt + Z_t^\top dW_t,\quad Y_T=\phi(X_T),
\end{equation}
where $Z_t=\sigma^\top\nabla_x V(t,X_t)$ is the gradient-weighted
diffusion process.  Han, Jentzen, and E~\cite{han2018solving} parameterize
$Z_t\approx\hat Z(X_t,t;\theta_t)$ by a neural network at each discretized
time step and minimize the terminal mismatch
$\mathbb{E}[|Y_T-\phi(X_T)|^2]$ via a forward simulation of the SDE.
The initial value $Y_0\approx\hat V(x_0,0)$ is recovered directly,
with no grid required.

Han and Long~\cite{han2020convergence} prove an \emph{a posteriori} error
bound: the error in $Y_0$ is bounded by the square root of the minimized
loss, provided the backward equation is decoupled from the forward
dynamics in its drift.
Negyesi, Huang, and Oosterlee~\cite{negyesi2024coupled} generalize this
to fully coupled FBSDEs.

\paragraph{Critical limitation: degenerate diffusion.}
The Feynman-Kac representation requires a non-degenerate diffusion in
\emph{all} state components.  If some state variables evolve
deterministically (i.e., rows of $\sigma$ are zero), the classical BSDE
formulation does not apply without a generalized extension.
This is a binding constraint for problems where part of the state---such
as an information or coverage variable---has no stochastic forcing.

\paragraph{Implementation note.}
The Julia script \texttt{deep\_bsde.jl} implements a
\emph{backward dynamic programming} (BDP) scheme rather than the full
Han et al.\ deep BSDE.  Starting from the terminal quadratic cost, it
fits a quadratic value model to the one-step Bellman backup:
\begin{equation}
  V_k(x) \approx \min_u\!\left[c(x,u)\Delta t
    + \mathbb{E}[V_{k+1}(x+\Delta x)]\right],
\end{equation}
exploiting the fact that the expectation is analytic for quadratic $V$
and additive Gaussian noise.  This is equivalent to the deep BSDE in the
LQR setting but does not generalize to non-quadratic value functions
without replacing the quadratic basis with a neural network and estimating
the expectation by Monte Carlo.

\paragraph{Key advantage.}
Unlike PINNs, the deep BSDE method requires only \emph{first-order}
gradients ($Z_t=\sigma^\top\nabla V$), avoiding the Hessian computation
needed for the Ito correction.  This makes it attractive for high-dimensional
problems where second-order AD is prohibitively expensive.
